\begin{helper}[The linear counting clause holds at unclamped windows]
  Let $k,n \in \mathbb{N}$, let $\delta,\epsilon,L \in \mathbb{R}$ and let
  $s \colon \mathbb{N} \to \mathbb{R}$. Assume:
  \begin{itemize}
    \item $s$ is monotone (non-decreasing) on $\set{0,1,\dots,k}$;
    \item $0 \le s(j) \le L$ for every $j \in \set{0,1,\dots,k}$;
    \item (clamped linear count) for all $a_0,b_0 \in \mathbb{R}$ with
      $0 \le a_0 \le b_0 \le L$ and $b_0 - a_0 \le 2\epsilon^{-1}\delta$,
      \[
        \#\set{ j \in \set{1,\dots,k} \colon a_0 \le s(j-1),\ s(j) \le b_0,\
          s(j)-s(j-1) < \delta } \le n
        .
      \]
  \end{itemize}
  Then for \emph{every} pair $a,b \in \mathbb{R}$ subject only to $b-a \le 2\epsilon^{-1}\delta$
  (in particular with no constraint $0 \le a \le b \le L$),
  \[
    \#\set{ j \in \set{1,\dots,k} \colon a \le s(j-1),\ s(j) \le b,\ s(j)-s(j-1) < \delta } \le n
    .
  \]

  \paragraph{Role.} The counting clause of the paper's $(\delta,\epsilon,n)$-curve condition
  (Definition delta-def, paper.tex 716--729) is stated only for windows $[a_0,b_0]$ that sit
  inside the parameter interval $[0,\length(\gamma)]$. Every consumer that pulls a window back
  along a reparametrisation --- \noderef{ArcDeltaEpsNCount} does this at three different
  $(a,b)$ pairs --- produces a window whose endpoints are shifted by the reparametrisation and
  therefore need not sit inside $[0,L]$ at all, even though its \emph{length} is still bounded by
  $2\epsilon^{-1}\delta$. The present statement removes that mismatch once and for all: the
  breakpoints of $s$ all lie in $[0,L]$, so the part of a window sticking out of $[0,L]$ is
  invisible to the count, and clamping the window to $[0,L]$ only shrinks its length.
\end{helper}
\begin{proof}
  Write
  \[
    F(p,q) := \set{ j \in \set{1,\dots,k} \colon p \le s(j-1),\ s(j) \le q,\ s(j)-s(j-1) < \delta }
  \]
  for $p,q \in \mathbb{R}$, and set $a_0 := \max(a,0)$, $b_0 := \min(b,L)$. We must show
  $\# F(a,b) \le n$.

  \emph{Case $F(a,b) = \emptyset$.} Then $\# F(a,b) = 0 \le n$.

  \emph{Case $F(a,b) \ne \emptyset$.} Fix $j_0 \in F(a,b)$, so $1 \le j_0 \le k$,
  $a \le s(j_0-1)$ and $s(j_0) \le b$. We verify the three hypotheses of the clamped linear count
  at the pair $(a_0,b_0)$.
  \begin{itemize}
    \item $0 \le a_0$: immediate from $a_0 = \max(a,0) \ge 0$.
    \item $a_0 \le b_0$: since $j_0 \ge 1$ we have $j_0 - 1 \le j_0 \le k$, so monotonicity of $s$
      gives $s(j_0-1) \le s(j_0)$. By hypothesis $a \le s(j_0-1)$, and $0 \le s(j_0-1)$ since
      $j_0-1 \le k$; hence $a_0 = \max(a,0) \le s(j_0-1)$. Likewise $s(j_0) \le b$ and
      $s(j_0) \le L$ since $j_0 \le k$; hence $s(j_0) \le \min(b,L) = b_0$. Chaining,
      $a_0 \le s(j_0-1) \le s(j_0) \le b_0$.
    \item $b_0 \le L$: immediate from $b_0 = \min(b,L) \le L$.
    \item $b_0 - a_0 \le 2\epsilon^{-1}\delta$: since $b_0 = \min(b,L) \le b$ and
      $a \le \max(a,0) = a_0$, we get $b_0 - a_0 \le b - a \le 2\epsilon^{-1}\delta$.
  \end{itemize}
  So the clamped linear count applies at $(a_0,b_0)$ and yields $\# F(a_0,b_0) \le n$.

  Finally $F(a,b) \subseteq F(a_0,b_0)$: let $j \in F(a,b)$, so $1 \le j \le k$, $a \le s(j-1)$,
  $s(j) \le b$ and $s(j)-s(j-1) < \delta$. Since $j-1 \le k$ we have $0 \le s(j-1)$, so
  $a_0 = \max(a,0) \le s(j-1)$; since $j \le k$ we have $s(j) \le L$, so
  $s(j) \le \min(b,L) = b_0$; and the edge condition is unchanged. Hence $j \in F(a_0,b_0)$.

  Therefore $\# F(a,b) \le \# F(a_0,b_0) \le n$.
\end{proof}
